\documentclass[11pt]{article}
\usepackage[margin=1.25in]{geometry}
\usepackage{natbib}
\usepackage{longtable}
\usepackage{booktabs}
\usepackage{subfig}
\usepackage{graphicx}

\begin{document}
\section{Results}
\subsection{Qualitative frequency estimates of cultural participation and culture talk in the NSS}
In the 2002 and 2004 NSS response categories for both culture talk and consumption frequencies for all five cultural domains are coded using qualitative frequency markers (e.g. ``very often,'' ``not very often,'' etc).  These types of items are also known in the literature as ``vague quantifiers.'' The use of qualitative frequency estimates has been thought of as problematic in the literature because they may introduce non-comparability across individuals.  For instance, one person's idea of the difference between ``very often'' and ``often'' may differ systematically from another person.\footnote{The only qualitative frequency indicator that appears to have an unambiguous meaning across respondents is, as should be expected, ``never.''}  For instance, research on cultural participation that compares frequency reports using vague quantifiers with reports by the same individuals expressed as self-reported \textit{counts} of the number of times that a person has engaged in the activity in the past year reveals a predictable bias:  socio-demographic groups who have a higher base propensity to engage in the activity more frequently (e.g. the young or the higher-educated) have a ``tougher'' (shifted to the left of the distribution of counts) threshold for what constitutes ``very often'' than those groups who are less likely to engage in the activity (as given by the self-reported counts).  

\subsection{Systematic differences in the matching of qualitative to quantitive frequencies across socio-demographic groups}
In the following analysis, the response variable is the individual's estimate of his or her own participation frequency expressed in terms of qualitative frequencies.  The main predictor is the same individual estimate of participation frequency but this time given by the estimated number of times that the person reports having engaged in the activity. My aim is to determine whether it is in fact the case that there exists structured heterogeneity in the way in which different persons think of the categories ``very often'' or ``fairly often.''  The guiding hypothesis is that this heterogeneity exists and that it is keyed to the formation of local norms among members of different socio-demographic groups.  In essence, when persons are asked to assign their own quantitative participation frequency into a qualitative slot, they take the modal participation frequency in their socio-demographic group as the ``norm.''  This means that for low participation groups, ``very often'' is matched to lower numeric frequencies than for high-participation groups and vice versa.

A simple ordered logit model of the association between qualitative and numeric estimates of participation frequency can be written as:

\begin{equation}
	logit[P([Y_q \leq j)]=\alpha_j + \phi Y_n	
	 \label{eq:equation1}
\end{equation}•

Where $Y_q$ are is the frequency of participation coded as qualitative frequency estimates and $\alpha_j$ are the ordered logit regression cutpoints.  Substantively, $\phi$ gives us the extent to which persons associate a given numerical frequency with a given qualitative frequency; if persons on average tend to match higher numeric frequencies to higher qualitative frequencies---as should be the obvious expectation---then the $\phi$ coefficient should be large and statistically significant.   

I test the hypothesis that different groups of people associate different numeric frequencies with different qualitative frequencies in a systematic way.  By systematic I mean that the association of quantity to quality is predictable by the extent to which different socio-demographic groups have different base expectations as to what constitutes frequent participation.  For instance, for socio-demographic groups for whom frequent cultural participation is the norm a given numeric frequency would be matched to a lower qualitative frequency in comparison to socio-demographic groups among whom frequent cultural participation is less common.  

I proceed by specifying the model:

\begin{equation}
	logit[P([Y_q \leq j)]=\alpha_j + \phi Y_n	+ \sum_k \beta_k X_k + \sum_k \gamma_k X_kY_n
	 \label{eq:equation2}
\end{equation}•

Here $X_k$ is a vector of individual socio-demographic characteristics, $\beta$ is a vector of coefficients that summarizes socio-demographic heterogeneity in participation frequency, and $\gamma$ is a vector of coefficients that summarize \emph{socio-demographic differences in the effect of the numeric indicator of frequency on the same variable expressed in qualitative terms}. Thus, if a given socio-demographic characteristic leads people to assign lower qualitative values to the same numeric response in comparison to other persons who are of low value in that characteristic, we should expect that $\gamma_k<0$. This would be the case for instance, if persons of low education tend to think of a given numeric frequency as meaning ``very often'' while persons of higher education tend to think of the same numeric frequency as meaning ``fairly often.'' The $\gamma$ interaction parameters are positive when being a given socio-demographic characteristic makes persons assign higher qualitative values to the same numeric frequency in comparison to persons who are low in that characteristic.  For instance, this would be the case if older individuals tend to think of a given numeric frequency as meaning ``fairly often'' (because they take as the norm the participation rates of their older peers) but younger people think of the same numeric participation frequency as actually meaning ``not very often.'' 

\subsection{Generalized Ordered Logit Model with Cutpoint-Shift}
The outcome items dealing with the frequency of culture talk are only available in the form of qualitative frequency estimates.  Given the results obtained above for the consumption items, it is safe to conclude that the individual's qualitative response is the observed realization of the person's unobserved ``objective'' (e.g. characterizable in numeric terms) frequency of culture-talk as filtered through the individual's propensity to assign his or her own numeric range to specific qualitative categories available in the survey.  I thus assume that all individual's have an ascertainable underlying propensity to engage in culture-talk, but that persons in different socio-demographic categories partition this continuum (e.g. assign specific ``cutpoints'' to the distribution) in systematically different ways.  

I use a generalization of the ordered logit model that allows for the indirect estimation of socio-demographic variation in cutpoint assignment in the absence of a numeric counterpart to the qualitatively coded response variable.  This model assumes that the observed qualitative response is the realization of an underlying continuous variable ($Y^*$)---representing the unobserved objective frequency----that is segmented at (in this case) four cutpoints ($\alpha_j$):

\begin{equation}
	y=g(Y^{*})=\left\{\begin{array}{rl}
	  1 &\mbox{if $Y^{*}p \leq \alpha_{1}$} \\
	  2 &\mbox{if $\alpha_{1} < Y^{*} \leq \alpha_{2}$}\\ 
	  3 &\mbox{if $\alpha_{2} <Y^{*} \leq \alpha_{3}$}\\ 
	  4 &\mbox{if $\alpha_{3} \leq Y^{*}$} 
	 \end{array} \right.
	 \label{eq:equation3}
\end{equation}
 
We only observe $Y=j$ if $\alpha_{j-1} < Y^{*} \leq \alpha_{j}$.  In the traditional ordered logit model, the $\alpha_j$ are a constant across all respondents.  That is, every person is assumed to map their underlying frequencies to the same qualitative responses.  By resorting to the generalized ordered logit model I relax this assumption, to accommodate for group-specific biases in the interpretation of the qualitative frequency categories. This model allows for $M-2$ of the cutpoints (where $M$ is the total number of ordered categories) to be a function of a selected set of socio-demographic variables:

\begin{equation}
\alpha_j=\hat{\alpha_j}+\gamma_kX'_k
\end{equation}•

For any given socio-demographic variable, the sign of the $\gamma_k$ coefficient tells us whether persons who are coded as high in the value of that variable tend to shift their cutpoints to the right or the left.  As shown in Figure~\ref{fig:cutpoint} the second cutpoint ($\alpha_2$) separates the ``not very often'' region from the ``fairly often'' segment of the distribution.  The third cutpoint ($\alpha_3$), in its turn, divides this last region from the ``very often'' part of the distribution.  Keeping in mind that cutpoints are reverse-coded (more negative cutpoints are assigned to ``higher'' categories of the outcome which indicate greater frequencies of culture talk), the substantive interpretation of the socio-demographic effects on the cutpoints is straighforward.  Aa negative effect on the cutpoint indicates that persons with high values on that predictor tend to shift that cutpoint to the \emph{right} of the underlying distribution; that is they apply more stringent standards to similar underlying numeric frequencies in order to categorize them as ``very often'' than persons low in that predictor.  A positive effect suggests the opposite; for individuals with high value on this predictor, the cutpoints are shifted to the \emph{left}; these persons judge lower underlying frequencies as sufficient for categorization as ``very often'' than persons with low values on this predictor.  Figure~\ref{fig:cutpoint} provides an intuitive illustration of the cutpoint shift phenomenon.
\newpage
{
\def\sym#1{\ifmmode^{#1}\else\(^{#1}\)\fi}
\begin{longtable}{l*{5}{c}}
\caption{Ordered logit models of qualitative frequency indicators: Interaction of socio-demographic effects with numeric indicators\label{tab:vague}}\\
\toprule\endfirsthead\midrule\endhead\midrule\endfoot\endlastfoot
                    &\multicolumn{1}{c}{Movies}&\multicolumn{1}{c}{Music}&\multicolumn{1}{c}{Books}&\multicolumn{1}{c}{Arts}&\multicolumn{1}{c}{Plays}\\
\midrule
main                &                   &                   &                   &                   &                   \\
Education           &       0.277       &       0.205       &       0.462\sym{*}&       0.704\sym{*}&       0.332\sym{*}\\
                    &      (1.51)       &      (0.71)       &      (2.00)       &      (3.33)       &      (1.90)       \\
Educ. X Numeric Freq.&      -0.134\sym{*}&     -0.0779       &      -0.177\sym{*}&      -0.358\sym{*}&      -0.142\sym{*}\\
                    &     (-1.91)       &     (-0.62)       &     (-1.93)       &     (-3.69)       &     (-1.91)       \\
Family Income       &       0.167       &      -0.142       &      -0.116       &      0.0534       &       0.277\sym{*}\\
                    &      (1.23)       &     (-0.67)       &     (-0.67)       &      (0.35)       &      (2.24)       \\
Income X Numeric Freq.&     -0.0325       &      0.0540       &      0.0295       &     0.00443       &     -0.0993\sym{*}\\
                    &     (-0.62)       &      (0.58)       &      (0.43)       &      (0.06)       &     (-1.85)       \\
Age                 &      -0.322\sym{*}&     -0.0575       &      -0.168       &    -0.00499       &      -0.235\sym{*}\\
                    &     (-2.84)       &     (-0.30)       &     (-1.30)       &     (-0.04)       &     (-2.31)       \\
Age X Numeric Freq. &      0.0830\sym{*}&     -0.0562       &       0.116\sym{*}&     0.00797       &      0.0981\sym{*}\\
                    &      (1.93)       &     (-0.69)       &      (2.25)       &      (0.15)       &      (2.29)       \\
Gender (Women=1)    &     -0.0400       &       0.371       &       0.182       &     -0.0387       &       0.771\sym{*}\\
                    &     (-0.09)       &      (0.55)       &      (0.36)       &     (-0.08)       &      (1.94)       \\
Women X Numeric Freq.&      0.0890       &     -0.0573       &      0.0369       &     -0.0238       &      -0.404\sym{*}\\
                    &      (0.54)       &     (-0.20)       &      (0.18)       &     (-0.11)       &     (-2.40)       \\
Race (Afr. Amer.=1) &       0.277       &      -2.458\sym{*}&      -0.272       &      -0.373       &      -0.544       \\
                    &      (0.40)       &     (-1.67)       &     (-0.31)       &     (-0.48)       &     (-0.84)       \\
Afr. Amer. X Numeric Freq.&       0.160       &       1.335\sym{*}&       0.291       &       0.664\sym{*}&       0.624\sym{*}\\
                    &      (0.61)       &      (1.91)       &      (0.82)       &      (1.76)       &      (2.17)       \\
Race (Asian=1)      &       2.127       &      -26.54       &      -2.250       &      -0.523       &      -1.102       \\
                    &      (1.37)       &     (-0.02)       &     (-0.93)       &     (-0.26)       &     (-0.62)       \\
Asian X Numeric Freq.&      -0.743       &       13.14       &       0.703       &       0.330       &      0.0842       \\
                    &     (-1.30)       &      (0.02)       &      (0.76)       &      (0.38)       &      (0.11)       \\
Race (Hispanic=1)   &       0.187       &      -0.427       &       0.776       &      -0.780       &      -0.594       \\
                    &      (0.18)       &     (-0.22)       &      (0.77)       &     (-0.65)       &     (-0.61)       \\
Hispanic X Numeric Freq.&       0.239       &       0.165       &     0.00837       &       0.500       &       0.489       \\
                    &      (0.61)       &      (0.21)       &      (0.02)       &      (0.95)       &      (1.20)       \\
Married             &       0.188       &      -0.352       &       0.228       &       0.185       &       0.787\sym{*}\\
                    &      (0.41)       &     (-0.47)       &      (0.41)       &      (0.36)       &      (1.80)       \\
Married X Numeric Freq.&      -0.122       &       0.158       &      -0.158       &   -0.000120       &      -0.400\sym{*}\\
                    &     (-0.69)       &      (0.49)       &     (-0.71)       &     (-0.00)       &     (-2.17)       \\
Urban Location      &       0.502       &      -0.564       &       0.199       &       0.539       &       0.571       \\
                    &      (1.13)       &     (-0.81)       &      (0.38)       &      (1.08)       &      (1.38)       \\
Urban X Numeric Freq.&      -0.238       &       0.258       &     -0.0137       &     -0.0698       &      -0.278       \\
                    &     (-1.39)       &      (0.87)       &     (-0.07)       &     (-0.31)       &     (-1.57)       \\
Movies Frequency (Numeric)&       4.736\sym{*}&                   &                   &                   &                   \\
                    &     (19.70)       &                   &                   &                   &                   \\
Music Frequency (Numeric)&                   &       2.494\sym{*}&                   &                   &                   \\
                    &                   &      (7.26)       &                   &                   &                   \\
Books Frequency (Numeric)&                   &                   &       3.801\sym{*}&                   &                   \\
                    &                   &                   &     (14.19)       &                   &                   \\
Arts Frequency (Numeric)&                   &                   &                   &       5.489\sym{*}&                   \\
                    &                   &                   &                   &     (18.88)       &                   \\
Plays Frequency (Numeric)&                   &                   &                   &                   &       5.161\sym{*}\\
                    &                   &                   &                   &                   &     (21.94)       \\
\midrule
cut1                &                   &                   &                   &                   &                   \\
Constant            &       7.491\sym{*}&       2.029       &       6.378\sym{*}&       11.95\sym{*}&       9.247\sym{*}\\
                    &      (7.43)       &      (1.21)       &      (5.37)       &     (10.20)       &      (9.94)       \\
\midrule
cut2                &                   &                   &                   &                   &                   \\
Constant            &       13.39\sym{*}&       4.147\sym{*}&       9.784\sym{*}&       17.21\sym{*}&       13.75\sym{*}\\
                    &     (12.44)       &      (2.47)       &      (7.97)       &     (14.05)       &     (14.21)       \\
\midrule
cut3                &                   &                   &                   &                   &                   \\
Constant            &       15.35\sym{*}&       5.746\sym{*}&       11.65\sym{*}&       20.62\sym{*}&       16.72\sym{*}\\
                    &     (14.18)       &      (3.42)       &      (9.40)       &     (16.39)       &     (16.85)       \\
\midrule
Observations        &        1786       &         995       &         989       &        1785       &        1784       \\
\bottomrule
\multicolumn{6}{l}{\footnotesize \textit{t} statistics in parentheses}\\
\multicolumn{6}{l}{\footnotesize \sym{*} \(p<0.10\)}\\
\end{longtable}
}

\newpage
{
\def\sym#1{\ifmmode^{#1}\else\(^{#1}\)\fi}
\begin{longtable}{l*{5}{c}}
\toprule\endfirsthead\midrule\endhead\midrule\endfoot\endlastfoot
                    &\multicolumn{1}{c}{Movies}&\multicolumn{1}{c}{Music}&\multicolumn{1}{c}{Books}&\multicolumn{1}{c}{Arts}&\multicolumn{1}{c}{Plays}\\
\midrule
Beta                &                   &                   &                   &                   &                   \\
Education           &       0.278\sym{*}&       0.363\sym{*}&       0.115\sym{*}&       0.118\sym{*}&       0.158\sym{*}\\
                    &      (4.89)       &      (5.28)       &      (2.07)       &      (2.73)       &      (3.79)       \\
Family Income       &      0.0358       &       0.110\sym{*}&      0.0689\sym{*}&       0.113\sym{*}&      0.0873\sym{*}\\
                    &      (1.22)       &      (2.78)       &      (1.69)       &      (3.24)       &      (2.48)       \\
Age                 &      -0.198\sym{*}&      -0.208\sym{*}&      -0.146\sym{*}&     -0.0708\sym{*}&     -0.0793\sym{*}\\
                    &     (-7.98)       &     (-4.81)       &     (-3.13)       &     (-2.58)       &     (-2.80)       \\
Gender (Women=1)    &       0.421\sym{*}&       0.205\sym{*}&       0.607\sym{*}&       0.318\sym{*}&       0.436\sym{*}\\
                    &      (4.70)       &      (1.74)       &      (5.00)       &      (3.32)       &      (4.72)       \\
Race (Afr. Amer.=1) &      -0.340       &      -0.406       &       0.541\sym{*}&       0.130       &       0.258       \\
                    &     (-1.51)       &     (-1.61)       &      (2.65)       &      (0.70)       &      (1.36)       \\
Race (Asian=1)      &    -0.00722       &     -0.0498       &      -0.134       &       0.121       &       0.257       \\
                    &     (-0.02)       &     (-0.10)       &     (-0.27)       &      (0.31)       &      (0.55)       \\
Race (Hispanic=1)   &      -0.884\sym{*}&     -0.0412       &      -0.165       &      -0.114       &      -0.287       \\
                    &     (-2.63)       &     (-0.16)       &     (-0.60)       &     (-0.43)       &     (-1.08)       \\
Married             &      0.0698       &      -0.218\sym{*}&     -0.0859       &     -0.0822       &     -0.0376       \\
                    &      (0.72)       &     (-1.67)       &     (-0.65)       &     (-0.78)       &     (-0.37)       \\
Urban Location      &      0.0641       &       0.301\sym{*}&       0.593\sym{*}&       0.146       &       0.283\sym{*}\\
                    &      (0.69)       &      (2.48)       &      (3.30)       &      (1.47)       &      (2.95)       \\
Movies Frequency (Numeric)&       1.468\sym{*}&                   &                   &                   &                   \\
                    &     (13.29)       &                   &                   &                   &                   \\
Music Frequency (Numeric)&                   &       0.546\sym{*}&                   &                   &                   \\
                    &                   &      (4.04)       &                   &                   &                   \\
Books Frequency (Numeric)&                   &                   &       1.457\sym{*}&                   &                   \\
                    &                   &                   &     (12.13)       &                   &                   \\
Arts Frequency (Numeric)&                   &                   &                   &       1.548\sym{*}&                   \\
                    &                   &                   &                   &     (20.72)       &                   \\
Plays Frequency (Numeric)&                   &                   &                   &                   &       1.169\sym{*}\\
                    &                   &                   &                   &                   &     (17.70)       \\
\midrule
Gamma\_2             &                   &                   &                   &                   &                   \\
Education           &      -0.132\sym{*}&      -0.160\sym{*}&                   &                   &                   \\
                    &     (-2.55)       &     (-2.68)       &                   &                   &                   \\
Race (Afr. Amer.=1) &       0.448\sym{*}&       0.776\sym{*}&                   &       0.497\sym{*}&       0.756\sym{*}\\
                    &      (2.36)       &      (3.94)       &                   &      (2.25)       &      (3.92)       \\
Race (Hispanic=1)   &       0.693\sym{*}&                   &                   &       0.758\sym{*}&       0.741\sym{*}\\
                    &      (2.50)       &                   &                   &      (2.50)       &      (2.51)       \\
Movies Frequency (Numeric)&      -0.481\sym{*}&                   &                   &                   &                   \\
                    &     (-4.67)       &                   &                   &                   &                   \\
Age                 &                   &       0.115\sym{*}&      0.0691       &      0.0789\sym{*}&       0.114\sym{*}\\
                    &                   &      (2.87)       &      (1.50)       &      (2.18)       &      (3.37)       \\
Music Frequency (Numeric)&                   &       0.281\sym{*}&                   &                   &                   \\
                    &                   &      (2.04)       &                   &                   &                   \\
Urban Location      &                   &                   &      -0.508\sym{*}&                   &                   \\
                    &                   &                   &     (-2.88)       &                   &                   \\
Books Frequency (Numeric)&                   &                   &      -0.177       &                   &                   \\
                    &                   &                   &     (-1.44)       &                   &                   \\
Family Income       &                   &                   &                   &     -0.0982\sym{*}&     0.00553       \\
                    &                   &                   &                   &     (-2.35)       &      (0.15)       \\
Race (Asian=1)      &                   &                   &                   &                   &       0.309       \\
                    &                   &                   &                   &                   &      (0.61)       \\
\midrule
Gamma\_3             &                   &                   &                   &                   &                   \\
Education           &      -0.246\sym{*}&      -0.233\sym{*}&                   &                   &                   \\
                    &     (-3.72)       &     (-2.65)       &                   &                   &                   \\
Race (Afr. Amer.=1) &       0.687\sym{*}&       0.828\sym{*}&                   &       0.485       &       0.486\sym{*}\\
                    &      (2.77)       &      (2.77)       &                   &      (1.47)       &      (1.69)       \\
Race (Hispanic=1)   &       1.110\sym{*}&                   &                   &       0.706       &       0.761\sym{*}\\
                    &      (3.21)       &                   &                   &      (1.57)       &      (1.81)       \\
Movies Frequency (Numeric)&      -0.712\sym{*}&                   &                   &                   &                   \\
                    &     (-5.85)       &                   &                   &                   &                   \\
Age                 &                   &      0.0651       &       0.132\sym{*}&      0.0183       &      0.0180       \\
                    &                   &      (1.23)       &      (2.44)       &      (0.34)       &      (0.37)       \\
Music Frequency (Numeric)&                   &       0.390\sym{*}&                   &                   &                   \\
                    &                   &      (1.80)       &                   &                   &                   \\
Urban Location      &                   &                   &      -0.479\sym{*}&                   &                   \\
                    &                   &                   &     (-2.27)       &                   &                   \\
Books Frequency (Numeric)&                   &                   &      -0.456\sym{*}&                   &                   \\
                    &                   &                   &     (-2.98)       &                   &                   \\
Family Income       &                   &                   &                   &      -0.180\sym{*}&     -0.0986\sym{*}\\
                    &                   &                   &                   &     (-2.92)       &     (-1.75)       \\
Race (Asian=1)      &                   &                   &                   &                   &       1.328\sym{*}\\
                    &                   &                   &                   &                   &      (2.24)       \\
\midrule
Alpha               &                   &                   &                   &                   &                   \\
\_cons\_1             &      -1.353\sym{*}&      -1.077\sym{*}&      -2.274\sym{*}&      -2.961\sym{*}&      -2.526\sym{*}\\
                    &     (-4.50)       &     (-2.13)       &     (-5.91)       &    (-12.02)       &    (-10.11)       \\
\_cons\_2             &      -1.605\sym{*}&      -3.246\sym{*}&      -3.611\sym{*}&      -5.154\sym{*}&      -5.092\sym{*}\\
                    &     (-6.09)       &     (-7.03)       &     (-9.85)       &    (-17.41)       &    (-18.19)       \\
\_cons\_3             &      -2.357\sym{*}&      -4.573\sym{*}&      -4.549\sym{*}&      -6.053\sym{*}&      -5.809\sym{*}\\
                    &     (-8.18)       &     (-6.92)       &    (-10.59)       &    (-16.55)       &    (-17.28)       \\
\midrule
Observations        &        1785       &         992       &         987       &        1783       &        1783       \\
\bottomrule
\multicolumn{6}{l}{\footnotesize \textit{t} statistics in parentheses}\\
\multicolumn{6}{l}{\footnotesize \sym{*} \(p<0.10\)}\\
\end{longtable}
}

\end{document}


\subsection{Differences in the marginal distribution of consumption versus talk across cultural domains}
Table~\ref{tab:gologit} shows the results of the generalized ordered logit estimation. As noted above, the binary indicator for outcome gives the us the difference in the marginal distribution between the consumption and talk outcomes.  A positive coefficient means that the distribution of talk is shifted to the \emph{right} in comparison to the distribution of consumption (more persons report talking about that cultural domain ``very often'' or ``fairly often'' than participating in it); a negative coefficient means that the distribution of culture talk is shifted to the \emph{left} in comparison to the consumption distribution (more persons report not engaging in culture talk in that domain in comparison to actually participating in it); a coefficient that is indistinguishable from zero tell us that talk and consumption are equitably distributed in that domain.  

The results show that for two cultural domains---book reading and live theater---the marginal distributions are statistically indistinguishable from one another.  For two other domains---film and the arts----talk is more prevalent than consumption.  Finally, in the case of music, consumption is more prevalent than talk.  This last result is consistent with research that shows music listening to be a relatively pervasive part of people's everyday life in relative independence from the social uses to which this consumption may be put.

\subsection{Socio-demographic heterogeneity in consumption versus talk}
The main effects of the socio-demographic predictors give us the coefficient estimates for the joint effect of that predictor on both consumption and culture-talk, under the assumption that this coefficient is the same for both outcomes in the matched pair design.  The \emph{interaction} between the matched-outcome binary indicator and the corresponding socio-demographic predictor is therefore a one-parameter test for whether that predictor has a statistically distinct effect (stronger if positive; weaker if negative) on talk in comparison to consumption \citep[Pp.]{agresti07}.  

Looking at the effects of education we find that across all five domains, the higher-educated are more likely to both engage in consumption and culture-talk across all five domains, with the education effects being generally stronger in those cultural domains associated with the non-profit fine arts.  In terms of the interaction between education and and the matched outcome I find that that culture-talk in the two commercial arts domains---film and music---is more strongly stratified by education than it is by consumption.  This is precisely the opposite of what I find in the case of the arts, which is a realm that appears to have stronger consumption stratification by education than it does in the case of talk ($t=-1.79$). This is different from what is observed in the case of book-reading and plays, where the null hypothesis that the education effect is the same across the two outcomes cannot be rejected within conventional levels of significance.  

Note that education is the only factor that displays this pattern; while wealthier persons tend to engage in more cultural talk and culture consumption across all five domains, these effects are generally weaker than the education effects and do not vary across the two outcomes.  In contrast to education, age has either null (in the case of books and arts) or negative effects on participation and talk (although the negative effect of age is relatively smaller for talk in relation to consumption).  Women consume \emph{and} talk about culture at higher rates than men, with the difference between men and women being much larger in the case of culture-talk than the corresponding effect in the film and theatrical domains.  In terms of the other socio-demographic predictors, I find that the effects of race and ethnicity tend to be weak and inconsistent, while being married depresses both talk and participation in the domains of literature, arts and theater; finally residing in an urban location increases both talk and participation and the effects are the same across the two outcomes.

\subsection{Socio-demographic heterogeneity in cutpoint location}
The last two sections of table~\ref{tab:gologit} shows the effect of selected socio-demographic covariates on the location of the second and third cutpoint of the underlying distribution.  

In the cutpoint equations, the effect of the binary indicator for the matched outcome is a single parameter test of the hypothesis that there exist distinct cutpoints for talk in relation to consumption within levels of all the predictors specified in the model.  This term essentially adjusts for the possibility that persons may interpret the distinction between the vague quantifiers ``fairly often'' and ``very often'' (for instance) differently in the case of talk than in the case of consumption. The results show that only movies appear to display a distinct pattern of cross-outcome shift in the cutpoints: essentially persons tend to have more relaxed standards as to what constitutes moderate and high frequency levels for movie-talk in comparison to actually \emph{going} to the movies, suggesting that frequent culture talk is more common (and therefore expected) in this domain than is sheer frequency of participation. This is to be expected since actually going out to the movies is a fairly stringent ``test'' of participation in this realm; persons are as likely to consume movies at home as they are to do so in the theater.  In the domain of music I find the opposite effect.  Here the talk cutpoint is shifted to the right, indicating that persons make a more demanding distinction as to what constitutes high-frequency of music-talk than they do of music-listening. This suggests that consumption without talk is more common and (statistically) normative in this domain than talk.

When it comes to socio-demographic heterogeneity in the placement of the cutpoints the results are generally as would be expected.  Education and income, both of which increase overall rates of talk and participation, shift the cutpoint location to the right across all five domains (with the partial exception of the arts).  This means that the higher educated and the wealthy make much more demanding judgments as to what constitutes fairly frequent versus very frequent behavior across both outcomes. For instance, two hypothetical respondents matched on every predictor but education would judge the same objective frequencies of talk and consumption differently, with the less educated being (probabilistically) more likely to deem as ``very often'' an objective frequency that the high-education respondent would classify as ``fairly often'' (the same goes for income).  

In the cutpoint equations, the interaction between the matched-outcome binary indicator and that predictor is---as in the main equation---a one-parameter test of the hypothesis that there exists structured socio-demographic heterogeneity in cutpoint location \emph{across} outcomes.  Generally, I find very limited evidence of this effect; instead, it is clear that most of the socio-demographic predictors have statistically indistinguishable effects on the cutpoint location across the two outcomes in all five cultural domains.  Two exceptions are notable.  First, education---as expected---shifts the cutpoint for book-reading \emph{consumption} rightwards---indicating that higher-educated members of the reading class have higher expectations as to what constitutes frequent reading than the lower educated---but shifts that of talk in that domain to the left; the same goes for income in relation to theater:  the wealthy have more demanding expectations as to what constitutes frequent theater going that the less wealthy but have opposite expectations when it comes to talk. 


